% =============================================================================
% Résumé (FR) et Abstract (EN)
% =============================================================================
\btpLiminary{Résumé}

Dans le secteur du Bâtiment et des Travaux Publics (BTP) au Cameroun, la
réponse aux Dossiers d'Appel d'Offres (DAO) constitue une activité stratégique
mais chronophage pour les bureaux d'études. Au sein de SITINFRA Cameroun, ce
processus demeure largement manuel~: l'analyse des pièces administratives,
techniques et financières ainsi que la préparation des offres mobilisent
plusieurs ingénieurs pendant de longues périodes, avec des risques d'omission
et une forte pression sur des délais compris entre 15 et 45~jours.

Ce mémoire pose la question de la conception d'un assistant intelligent capable
d'exploiter de manière fiable et sécurisée le patrimoine documentaire de
l'entreprise afin d'automatiser l'analyse des DAO et la génération de réponses
conformes au Code des Marchés Publics camerounais. L'approche retenue repose
sur la Génération Augmentée par Récupération (\anglais{Retrieval-Augmented
Generation}, RAG), enrichie d'une ingestion contextuelle, d'une récupération
hybride, d'un mécanisme correctif (CRAG) et d'une orchestration maîtrisée, afin
de concilier la performance des grands modèles de langage avec les exigences de
conformité, de traçabilité et de sécurité documentaire.

Une démarche méthodologique fondée sur la science de la conception
(\anglais{Design Science Research}) structure la modélisation de la solution, et
un protocole d'évaluation moderne (RAGAS, \anglais{LLM-as-a-Judge}) est proposé
pour valider les hypothèses de recherche. Les résultats obtenus établissent une
base technique robuste pour l'assistance à la préparation des offres et
l'amélioration de la compétitivité des bureaux d'études du BTP.

\vspace{0.6em}
\begin{btpbox}[Mots-clés]
Intelligence artificielle~; Génération Augmentée par Récupération (RAG)~;
Dossier d'Appel d'Offres (DAO)~; BTP~; marchés publics~; capitalisation
documentaire~; supervision humaine~; SITINFRA Cameroun.
\end{btpbox}

\btpEndMark
\clearpage

% -----------------------------------------------------------------------------
\btpLiminary{Abstract}
\selectlanguage{english}

In Cameroon's Building and Public Works (BTP) sector, responding to tender
documents (\anglais{Dossiers d'Appel d'Offres}, DAO) is a strategic yet
time-consuming activity for engineering consultancies. At SITINFRA Cameroun,
this process remains largely manual: analysing administrative, technical and
financial documents and preparing bids mobilises several engineers over long
periods, with a high risk of omissions and strong pressure within deadlines of
15 to 45 days.

This thesis addresses the design of an intelligent assistant able to exploit the
company's documentary assets in a reliable and secure way, in order to automate
the analysis of tender documents and to generate responses that comply with the
Cameroonian Public Procurement Code. The proposed approach is based on
Retrieval-Augmented Generation (RAG), enhanced with contextual ingestion,
hybrid retrieval, a corrective mechanism (CRAG) and controlled orchestration,
so as to reconcile the performance of large language models with the
requirements of compliance, traceability and documentary security.

A Design Science Research methodology structures the modelling of the solution,
and a modern evaluation protocol (RAGAS, LLM-as-a-Judge) is proposed to validate
the research hypotheses. The results establish a robust technical foundation for
supporting bid preparation and improving the competitiveness of BTP engineering
firms.

\vspace{0.6em}
\begin{btpbox}[Keywords]
Artificial Intelligence; Retrieval-Augmented Generation (RAG); Tender Documents
(DAO); Public Works (BTP); Public Procurement; documentary capitalisation;
human-in-the-loop; SITINFRA Cameroun.
\end{btpbox}

\selectlanguage{french}
\btpEndMark
\clearpage
